% CLAUDE REPORT::BRIDGE-0001::2026-11-12T20:55Z
% !TEX program = xelatex
% Compiles with XeLaTeX or pdfLaTeX (ASCII source).
\documentclass[11pt]{article}
\usepackage[margin=1in]{geometry}
\usepackage{amsmath,amssymb,amsthm}
\usepackage{xurl}
\newtheorem{theorem}{Theorem}
\newtheorem{proposition}[theorem]{Proposition}
\newtheorem{corollary}[theorem]{Corollary}
\theoremstyle{definition}
\newtheorem{definition}{Definition}
\theoremstyle{remark}
\newtheorem*{remark}{Remark}
\newcommand{\SAT}{\mathrm{SAT}}
\newcommand{\Wb}{\overline W}
\newcommand{\status}[1]{\textbf{[#1]}}

\title{CLAUDE BRIDGE-0001\\[2pt]
\large Where the weakest bridge $\SAT\in P\Rightarrow \Wb_R(N)\le\mathrm{poly}(N)$ is\\
a reformulation, and where it is not}
\author{CLAUDE (scheduled routine; supplementary, not an LSPV role)}
\date{2026-11-12 12:55 America/Los\_Angeles}

\begin{document}
\maketitle

\noindent\textbf{Status.} P versus NP remains \textbf{OPEN}. This report proves elementary
classification facts (\status{PROVED/EXACT}) about the shared worst-case target of the
2026-11-02 refocus and records known theorems (\status{KNOWN THEOREM}) that bear on it. It proves no
lower bound and no bridge. It is additive: no schedules, allocations, ownership regions or root
\texttt{main.tex} were changed (this routine owns no \texttt{main.tex} region).

\section{Setting}
Fix one sound, complete, effectively specified reduction system $R$ (for instance the project's
mixed calculus $R_{\mathrm{mix341}}$, or the multisign DAG-hypergraph reduction calculus) with
nonnegative node weights. As in
\texttt{research/coordination/2026-11-02\_1204\_WORST\_CASE\_MINIMUM\_DERIVATION\_FIVE\_ROLE\_REFOCUS.md},
\[
W_R(F)=\min_{\pi:F\vdash_R\bot}\operatorname{cost}_R(\pi),\qquad
\Wb_R(N)=\max_{F\in\mathrm{UNSAT},\ |\langle F\rangle|\le N} W_R(F).
\]
The inner minimum grants the best reduction sequence for free. This is a legitimate concession for a
\emph{lower bound}: any procedure that actually executes an $R$-derivation pays at least $W_R(F)$,
and searching for the sequence only adds work.

\begin{definition}
$\mathrm{LB}(R)$: for every $k$ and $N_0$ there is $N\ge N_0$ with $\Wb_R(N)>(N+1)^k$.
$\mathrm{PB}(R)=\neg\mathrm{LB}(R)$ (cost-polynomially bounded).
$\mathrm{B3}(R)$: $\SAT\in P\Rightarrow \mathrm{PB}(R)$.
$R'$ \emph{cost-simulates} $R$ if there is a polynomial $q$ with
$W_{R'}(F)\le q(|\langle F\rangle|+W_R(F))$ for all UNSAT $F$.
\end{definition}

The program's plan is $\mathrm{LB}(R)\wedge\mathrm{B3}(R)\Rightarrow \SAT\notin P$. Both halves are
needed for the \emph{same} $R$.

\section{Results}

\begin{theorem}[Squeeze]\status{PROVED/EXACT}\label{thm:squeeze}
If $\mathrm{LB}(R)$ holds, then $\mathrm{B3}(R)\iff \SAT\notin P$.
\end{theorem}
\begin{proof}
$\mathrm{B3}(R)$ is the disjunction $\SAT\notin P \vee \mathrm{PB}(R)$, and $\mathrm{PB}(R)$ is false.
\end{proof}

\begin{corollary}\status{PROVED/EXACT given the cited KNOWN THEOREMS}
For tree and DAG Resolution (Haken 1985; Urquhart 1987), Cutting Planes (Pudl\'ak 1997),
Polynomial Calculus (Razborov 1998) and bounded-depth Frege (Ajtai 1988; Kraj\'i\v{c}ek--Pudl\'ak--Woods,
Pitassi--Beame--Impagliazzo 1995), with unit or polynomially related node weights, the
``weakest bridge'' $\mathrm{B3}(R)$ is \emph{literally equivalent to} $P\ne NP$.
Proving it is exactly as hard as the whole problem.
If the reported AC$^0[p]$-Frege lower bound (Lu--Santhanam--Tzameret, ECCC TR25-134) holds as
stated, the same applies to AC$^0[p]$-Frege (\status{CANDIDATE}, not re-verified here).
\end{corollary}

\begin{theorem}[Monotonicity]\status{PROVED/EXACT}\label{thm:mono}
If $R'$ cost-simulates $R$, then $\mathrm{PB}(R)\Rightarrow\mathrm{PB}(R')$,
$\mathrm{LB}(R')\Rightarrow\mathrm{LB}(R)$, and $\mathrm{B3}(R)\Rightarrow\mathrm{B3}(R')$.
\end{theorem}
\begin{proof}
If $\Wb_R(N)\le p(N)$ then $\Wb_{R'}(N)\le q(N+p(N))$. The other two statements follow by
contraposition and by composing with the hypothesis of B3.
\end{proof}

So making $R$ stronger makes the bridge easier and the lower bound harder. A route that is not a
reformulation must choose $R$ in the band where $\mathrm{LB}(R)$ is \emph{not} already known. By
Theorem~\ref{thm:squeeze}, this is a necessary condition. Mixing a lower-bounded subsystem into a
stronger $R$ does not move $R$ out of that band, because the stronger system can use either part.

\begin{theorem}[World classification]\status{PROVED/EXACT}\label{thm:worlds}
In every world with $P\ne NP$ (Heuristica, Pessiland, Minicrypt, Cryptomania), $\mathrm{B3}(R)$ is
true for every $R$, vacuously. In Algorithmica, $\mathrm{B3}(R)\iff\mathrm{PB}(R)$. Hence:
\begin{itemize}
\item if $\mathrm{LB}(R)$ is a theorem, B3$(R)$ is ``Equivalent to $P\ne NP$'' (ledger row L41);
\item for $R=$ Extended Frege, B3 is open in Algorithmica, so it is a \emph{half-win}
(ledger row L42): proving it outright is strictly less than proving $P\ne NP$, and refuting it
proves $P=NP$.
\end{itemize}
\end{theorem}

\begin{proposition}[Provability upgrade]\status{KNOWN THEOREM, restated}\label{prop:pv}
Kraj\'i\v{c}ek: PV proves $P=NP$ if and only if PV defines a polynomial-time function that outputs,
for each formula, either a satisfying assignment or an Extended Frege refutation. With Cook's 1975
translation, a PV proof of such correctness yields polynomial-size EF refutations of every UNSAT
formula. Consequently
\[
H_{\mathrm{PV}}:\quad \SAT\in P\ \Rightarrow\ \text{PV proves } P=NP \text{ for some PV function}
\]
implies $\mathrm{B3}(\mathrm{EF})$. The same holds for any sound theory $T\supseteq\mathrm{PV}$ and its
reflection system $R_T$, whose refutations of $F$ are pairs (a $T$-proof that a clocked algorithm $g$
is sound on the UNSAT side, a run of $g$ on $F$). Soundness of $R_T$ needs the $\Pi_1$-soundness of $T$.
\end{proposition}

\begin{remark}[What the gap actually is]
Under $\SAT\in P$ with solver $f$, the SAT side certifies itself: search-to-decision gives an
assignment that can be checked. The UNSAT side does not. A rejecting run of $f$ is an $R_T$-proof only if
$T$ proves that $f$ never wrongly answers UNSAT. So for strong $R$, the full content of B3 is the
\emph{provability gap} $H_T$: ``fast correct'' does not imply ``fast and provably correct.'' This
matches the corpus's trace-correctness firewall (LSPV-5 REFEREE294) and the generator barrier
(LSPV-1/335 Thm~2, LSPV-1/334). Raising $T$ (PV $\to S^1_2\to$ PA $\to$ ZFC) weakens $H_T$ and
strengthens $R_T$. By Theorem~\ref{thm:mono}, it also makes $\mathrm{LB}(R_T)$ harder.
Moreover $\mathrm{LB}(R_T)$ alone already implies that no polynomial-time algorithm is both
$T$-provably UNSAT-sound and actually complete. That is a consistency statement, not a separation.
\end{remark}

\begin{corollary}[Consequence for the five-role program]\status{PROVED/EXACT, strategic}
(a) A superpolynomial $\Wb_R$ lower bound for a system at or below the known-lower-bound frontier,
combined with B3 for that same system, is a restatement of $P\ne NP$. Hard UNSAT families for
Resolution-like fragments of $R_{\mathrm{mix341}}$ (Tseitin, PHP, the multisign hypergraph
families) are useful controls but cannot carry the bridge.
(b) The only regime where the decomposition can help is a system $R$ that cost-simulates EF (or is
not known to be lower-bounded), together with a proof of $H_T$ or of an equivalent provability
upgrade for that $R$. Both obligations are \status{OPEN}. A superpolynomial EF lower bound is
itself a major open problem (Pich--Santhanam 2023 obtain $P\ne NP$ from it only under additional
hypotheses).
(c) For the project's $R_{\mathrm{mix341}}$, the decisive question is therefore: does $R_{\mathrm{mix341}}$
cost-simulate EF? If no, $\mathrm{B3}(R_{\mathrm{mix341}})$ has no known route other than proving $P\ne NP$ directly.
If yes, its lower bound is at least as hard as an EF lower bound. Either answer is useful, and this
should be the next audit item for the Time-Space Resource Auditor and the Red Team.
\end{corollary}

\section{Literature checked this run}
\begin{itemize}
\item Kraj\'i\v{c}ek, notes on PV and $P=NP$ (equivalence used in Prop.~\ref{prop:pv}):
\url{https://webarchiv.lib.cas.cz/wayback/20090227233806/http:/www.math.cas.cz/~krajicek/steve.pdf}.
\item Pich--Santhanam, \emph{Towards P$\ne$NP from Extended Frege lower bounds}, ECCC TR23-199,
\url{https://eccc.weizmann.ac.il/report/2023/199}; arXiv 2312.08163.
\item Atserias et al., \emph{The Proof Analysis Problem} (2025), arXiv 2506.16956: PV1 proves
correctness of a specific extraction algorithm with short EF proofs via Cook's translation.
This is an exact instance of the Prop.~\ref{prop:pv} mechanism, for a restricted problem only.
\item Lu--Santhanam--Tzameret, ECCC TR25-134 (AC$^0[p]$-Frege lower bound for DNFs; reported, not re-verified).
\item Cook--Reckhow 1979; Haken 1985; Urquhart 1987; Pudl\'ak 1997; Razborov 1998; Ajtai 1988 (classical).
\end{itemize}
Barrier-transfer: the closest structural analogue outside complexity is \emph{verified compilation}
(software verification). A correct but unverified optimizer does not yield proof-carrying code.
Classification: useful structural analogy; the exact missing transfer theorem is $H_T$.
Post-quantum ledger: no item applies. Item 6 (special soundness) is the nearest analogue, but it is
only an ANALOGY, because cryptographic binding assumptions are inadmissible.

\section{Required audit entries}
\textbf{CORPUS REVIEWED.} AGENTS.md (including the 2026-11-09 knot/link handoff header); the
2026-11-02 worst-case refocus; LSPV-1/335 (core reduction, generator barrier); LSPV-1/367
(reflection inherits Resolution hardness); LSPV-1/539--540 sections and the status header of root
\texttt{main.tex}; the coordination handoffs of 2026-11-12 (LSPV-1/598--599, LSPV-2/581,
LSPV-3/594--595, LSPV-4/567, LSPV-5/509); repository search for Kraj\'i\v{c}ek, PV, Haken, B3 and
reflection (92 files mention Kraj\'i\v{c}ek); the CLAUDE world ledger and its prior reports.

\textbf{DUPLICATION / CONTRADICTION AUDIT.} The corpus already records that the parity recognizer
refutes B1 for Resolution, the generator barrier (LSPV-1/335), and the trace-correctness firewall
(REFEREE294). I found no corpus statement of the squeeze (Thm~\ref{thm:squeeze}) or of the
monotone tradeoff (Thm~\ref{thm:mono}) in this form. Both are elementary, so I claim no
mathematical novelty, only that they are now stated explicitly. I found no contradictions.

\textbf{WHAT CHANGED RELATIVE TO ROOT MAIN.} Nothing in \texttt{main.tex}. New world-ledger rows
L41--L45.

\textbf{ADOPTED / REJECTED / RETIRED PRIOR COMPONENTS.} Adopted: the refocus definitions of
$\Wb_R$ and B1/B2/B3. Rejected as a bridge carrier: B3 for any lower-bounded fragment, which is a
reformulation (not retired as a control).

\textbf{STALE OR OVERSTATED CLAIMS FOUND.} None found. Any future text calling B3 ``weaker than
P$\ne$NP'' should name the system: for lower-bounded systems, B3 is equivalent to P$\ne$NP.

\textbf{Predecessor check.} This routine is outside the LSPV stagger. The newest LSPV artifacts
(LSPV-4/568, commits up to \texttt{bdb222dc}) exist on the branch. No recovery action was taken.

\textbf{Artifacts.} Ledger builder \texttt{research/agents/claude/world-ledger/build\_ledger.py}
(rows L41--L45, regenerated CSV and XLSX; the auto-categories were Equivalent-to-P$\ne$NP,
Half-win, Half-win, Theorem and Theorem, all with consistency OK). Living log:
\texttt{research/agents/claude/LIVING\_LOG.md}. The build run takes under 2\,s, so no chunking was
needed. No LaTeX engine is installed in this container; structural checks (balanced braces and environments) passed, but the file was not compiled.

\noindent P versus NP remains \textbf{OPEN}.
\end{document}
% END_OF_VERIFIED_REPORT::CLAUDE::BRIDGE-0001
